%! Parametric Functions and Rates of Change — Unit 8, Lesson 8.2 — Guided Notes
\documentclass[10pt]{article}
\usepackage{precalculus-article}
\usepackage{precalculus-boxes}
\newcommand{\TallMath}[1]{$\displaystyle #1\rule[-1.4em]{0pt}{3.2em}$}

\begin{document}

\pageheader{Unit 8, Lesson 8.2}{Guided Notes}
\namedateperiod

\vspace{0.10in}

%----------------------------------------------------------------------
% Objectives
%----------------------------------------------------------------------
\begin{objectivebox}
\textbf{I will be able to\ldots}
\begin{enumerate}[leftmargin=1.5em, itemsep=4pt]
  \item \textbf{LO 4.3.A} --- Identify the direction of planar motion
    at any value of the parameter by determining whether $x(t)$ and
    $y(t)$ are increasing or decreasing at that moment. \writeline
  \item \textbf{LO 4.3.A} --- Explain why the same curve can be traversed
    in different directions by different parametrizations. \writeline
  \item \textbf{LO 4.3.A} --- Compute the average rate of change of $x(t)$
    and $y(t)$ over $[t_1, t_2]$ and use their ratio to find the slope
    of the parametric curve between two points. \writeline
\end{enumerate}
\end{objectivebox}

\vspace{0.08in}

%----------------------------------------------------------------------
% Vocabulary
%----------------------------------------------------------------------
\begin{vocabbox}
\begin{description}[leftmargin=0pt, itemsep=6pt]
  \item[\termblanklong{direction of motion}]
    Determined by whether each component is
    \underline{\hspace{2.5cm}} or \underline{\hspace{2.5cm}}.
    If $x(t)$ is increasing $\Rightarrow$ moving \underline{\hspace{1.5cm}};
    if decreasing $\Rightarrow$ moving \underline{\hspace{1.5cm}}.
    If $y(t)$ is increasing $\Rightarrow$ moving \underline{\hspace{1.5cm}};
    if decreasing $\Rightarrow$ moving \underline{\hspace{1.5cm}}.

  \item[\termblanklong{ARC of $x(t)$}]
    The average rate of change of $x(t)$ over $[t_1, t_2]$:
    \quad $\dfrac{x(t_2) - x(t_1)}{t_2 - t_1}$

  \item[\termblanklong{ARC of $y(t)$}]
    The average rate of change of $y(t)$ over $[t_1, t_2]$:
    \quad $\dfrac{y(t_2) - y(t_1)}{t_2 - t_1}$

  \item[\termblanklong{slope of parametric curve}]
    The slope of the \underline{\hspace{2cm}} line connecting two points
    on the curve:
    \quad $\dfrac{\Delta y}{\Delta x} = \dfrac{\text{ARC of } y(t)}{\text{ARC of } x(t)}$
    \quad provided ARC of $x(t) \ne 0$.

  \item[\termblanklong{reparametrization}]
    A different pair of component functions that traces the
    \underline{\hspace{2.5cm}} set of points in the plane, possibly in a
    \underline{\hspace{2cm}} direction.
\end{description}
\end{vocabbox}

\vspace{0.08in}

%----------------------------------------------------------------------
% Hook
%----------------------------------------------------------------------
\begin{hookbox}
\textbf{Entry Question:} Two drones follow the same figure-eight path ---
Drone~A clockwise, Drone~B counterclockwise.

\vspace{4pt}
\textbf{Q1.} Can two different parametric functions trace the same path?
What would have to be different about them?
\writelines{2}

\textbf{Q2.} How do you determine which direction a particle is moving
horizontally at a specific value of $t$?
\writeline
\end{hookbox}

\vspace{0.08in}

%----------------------------------------------------------------------
% Step 1: Direction of Motion
%----------------------------------------------------------------------
\begin{notesbox}{LO 4.3.A --- Direction of Planar Motion (EK 4.3.A.1--A.2)}

\textbf{Key rule:} Analyze $x(t)$ and $y(t)$ \emph{independently}.

\begin{center}
\renewcommand{\arraystretch}{1.5}
\begin{tabular}{|l|l|}
\hline
If $x(t)$ is \textbf{increasing} & particle moves \blank{2.5cm} \\
\hline
If $x(t)$ is \textbf{decreasing} & particle moves \blank{2.5cm} \\
\hline
If $y(t)$ is \textbf{increasing} & particle moves \blank{2.5cm} \\
\hline
If $y(t)$ is \textbf{decreasing} & particle moves \blank{2.5cm} \\
\hline
\end{tabular}
\end{center}

\vspace{8pt}
\textbf{Example:} Let $f(t) = (t^2 - 2t,\; 3t - t^2)$ for $0 \le t \le 4$.

\vspace{6pt}
\textbf{Step 1 --- Build a table of values.}

\vspace{4pt}
\begin{center}
\renewcommand{\arraystretch}{1.7}
\begin{tabular}{|c|c|c|c|}
\hline
$t$ & $x(t) = t^2 - 2t$ & $y(t) = 3t - t^2$ & $(x,\; y)$ \\
\hline
$0$ & \blank{2cm} & \blank{2cm} & \blank{2.8cm} \\
\hline
$1$ & \blank{2cm} & \blank{2cm} & \blank{2.8cm} \\
\hline
$2$ & \blank{2cm} & \blank{2cm} & \blank{2.8cm} \\
\hline
$3$ & \blank{2cm} & \blank{2cm} & \blank{2.8cm} \\
\hline
$4$ & \blank{2cm} & \blank{2cm} & \blank{2.8cm} \\
\hline
\end{tabular}
\end{center}

\vspace{8pt}
\textbf{Step 2 --- Direction at $t = 0.5$.}

Compare $x(0)$ and $x(1)$: $x(0) = $ \blank{1.2cm}, $x(1) = $ \blank{1.2cm}.

Since $x(1)$ \blank{1.2cm} $x(0)$, the function $x(t)$ is \blank{2cm} near $t = 0.5$.

$\Rightarrow$ particle is moving \blank{2.2cm}

Compare $y(0)$ and $y(1)$: $y(0) = $ \blank{1.2cm}, $y(1) = $ \blank{1.2cm}.

Since $y(1)$ \blank{1.2cm} $y(0)$, the function $y(t)$ is \blank{2cm} near $t = 0.5$.

$\Rightarrow$ particle is moving \blank{2.2cm}

\vspace{8pt}
\textbf{Step 3 --- Direction at $t = 2.5$.}

Compare $x(2)$ and $x(3)$: is $x$ increasing or decreasing? \blank{2.5cm}

$\Rightarrow$ particle is moving \blank{2.2cm}

Compare $y(2)$ and $y(3)$: is $y$ increasing or decreasing? \blank{2.5cm}

$\Rightarrow$ particle is moving \blank{2.2cm}

\end{notesbox}

\vspace{0.08in}

%----------------------------------------------------------------------
% Step 2: Same Curve, Different Directions
%----------------------------------------------------------------------
\begin{notesbox}{LO 4.3.A --- Same Curve, Different Parametrizations (EK 4.3.A.3)}

\textbf{Key idea:} The same set of points in the plane can be traced by
different parametric functions --- possibly in different directions.

\vspace{6pt}
\textbf{Example:} Let $f(t) = (t,\; t^2)$ and $g(t) = (3 - t,\; (3-t)^2)$
for $0 \le t \le 3$.

\vspace{4pt}
\begin{center}
\renewcommand{\arraystretch}{1.7}
\begin{tabular}{|c|c|c||c|c|}
\hline
& \multicolumn{2}{c||}{Robot A: $f(t)=(t,\;t^2)$}
& \multicolumn{2}{c|}{Robot B: $g(t)=(3{-}t,\;(3{-}t)^2)$} \\
\hline
$t$ & $(x,\;y)$ for $f$ & direction & $(x,\;y)$ for $g$ & direction \\
\hline
$0$ & \blank{2.2cm} & & \blank{2.2cm} & \\
\hline
$1$ & \blank{2.2cm} & & \blank{2.2cm} & \\
\hline
$2$ & \blank{2.2cm} & & \blank{2.2cm} & \\
\hline
$3$ & \blank{2.2cm} & & \blank{2.2cm} & \\
\hline
\end{tabular}
\end{center}

\vspace{6pt}
Do $f$ and $g$ trace the same set of points? \blank{1.5cm}

Are they moving in the same direction? \blank{1.5cm}

\end{notesbox}

\vspace{0.08in}

%----------------------------------------------------------------------
% Step 3: Slope via ARC Ratio
%----------------------------------------------------------------------
\begin{notesbox}{LO 4.3.A --- Slope of the Curve via Average Rates of Change (EK 4.3.A.4)}

\textbf{Formula:} Over the interval $[t_1, t_2]$:
\[
  \text{slope of curve} = \frac{\Delta y}{\Delta x}
  = \frac{y(t_2) - y(t_1)}{x(t_2) - x(t_1)}
  \quad \text{(provided } x(t_2) \ne x(t_1)\text{)}
\]

\vspace{6pt}
\textbf{Example (continued):} Use $f(t) = (t^2 - 2t,\; 3t - t^2)$.
Find the slope of the curve from $t = 1$ to $t = 3$.

\vspace{6pt}
$x(1) = $ \blank{1.5cm} \qquad $x(3) = $ \blank{1.5cm} \qquad $\Delta x = $ \blank{1.5cm}

$y(1) = $ \blank{1.5cm} \qquad $y(3) = $ \blank{1.5cm} \qquad $\Delta y = $ \blank{1.5cm}

\vspace{6pt}
Slope $= \dfrac{\Delta y}{\Delta x} = \dfrac{\blank{1.5cm}}{\blank{1.5cm}} = $ \blank{1.5cm}

\vspace{6pt}
\textbf{Interpretation:} The secant line connecting the points on the
curve at $t = 1$ and $t = 3$ has slope \blank{1.5cm}.

\end{notesbox}

\end{document}
